% ============================================================
% CHAPITRE 2 : Transformée de Laplace & Fonctions de transfert
% ============================================================
\chapterbanner{Chapitre 2 --- Laplace \& Fonctions de transfert}{ch2orange}

\begin{tcolorbox}[colback=ch2orange!4,colframe=ch2orange!80!black,title={Table de Laplace essentielle}]
\begin{center}
\renewcommand{\arraystretch}{1.05}
\scriptsize
\begin{tabular}{@{}lll@{}}
\toprule
$f(t),\; t \geq 0$ & $F(s) = \Lap\{f(t)\}$ & Note \\
\midrule
$\delta(t)$ & $1$ & Impulsion \\
$\mathbf{1}(t)$ & $1/s$ & \'Echelon \\
$t$ & $1/s^2$ & Rampe \\
$t^n/n!$ & $1/s^{n+1}$ & \\
$e^{-at}$ & $1/(s{+}a)$ & Exponentielle \\
$t\,e^{-at}$ & $1/(s{+}a)^2$ & \\
$\frac{t^{n-1}}{(n{-}1)!}e^{-at}$ & $1/(s{+}a)^n$ & P\^ole multiple \\
$\sin(\omega t)$ & $\omega/(s^2{+}\omega^2)$ & \\
$\cos(\omega t)$ & $s/(s^2{+}\omega^2)$ & \\
$e^{-at}\sin(\omega t)$ & $\omega/[(s{+}a)^2{+}\omega^2]$ & \\
$e^{-at}\cos(\omega t)$ & $(s{+}a)/[(s{+}a)^2{+}\omega^2]$ & \\
$1 - e^{-t/\tau}$ & $1/[s(1{+}\tau s)]$ & R\'ep.ind. 1\textsuperscript{er} o. \\
\bottomrule
\end{tabular}
\end{center}
\end{tcolorbox}

\begin{tcolorbox}[colback=ch2orange!2,colframe=ch2orange!80!black,title={Propri\'et\'es de la transform\'ee}]
\begin{center}
\renewcommand{\arraystretch}{1.05}
\scriptsize
\begin{tabular}{@{}ll@{}}
\toprule
Propri\'et\'e & Formule \\
\midrule
Lin\'earit\'e & $\Lap\{af{+}bg\} = aF{+}bG$ \\
D\'erivation & $\Lap\{\dot{f}\} = sF(s) - f(0^-)$ \\
D\'eriv\'ee $n$-i\`eme & $\Lap\{f^{(n)}\} = s^n F - \sum_{k=0}^{n-1} s^{n-1-k} f^{(k)}(0^-)$ \\
Int\'egration & $\Lap\{\int_0^t f\,d\tau\} = F(s)/s$ \\
Retard temporel & $\Lap\{f(t{-}T_r)\cdot\mathbf{1}(t{-}T_r)\} = e^{-sT_r}F(s)$ \\
D\'ecalage fr\'eq. & $\Lap\{e^{-at}f(t)\} = F(s{+}a)$ \\
Val. initiale & $\lim_{t\to 0^+} f(t) = \lim_{s\to\infty} sF(s)$ \\
\textbf{Val. finale} & $\boldsymbol{\lim_{t\to\infty} f(t) = \lim_{s\to 0} sF(s)}$ \\
Convolution & $\Lap\{f*g\} = F(s)\cdot G(s)$ \\
\bottomrule
\end{tabular}
\end{center}
\piege{Th. valeur finale : valable \textbf{ssi} tous les p\^oles de $sF(s)$ ont Re~$< 0$ (sauf \'event. un p\^ole simple en $s{=}0$). Ne pas appliquer si syst\`eme instable !}
\end{tcolorbox}

\begin{tcolorbox}[colback=ch2orange!4,colframe=ch2orange!80!black,title={Fonction de transfert}]
\tightmath{
\boxed{G(s) = \frac{Y(s)}{U(s)} = \frac{b_m s^m + \cdots + b_0}{a_n s^n + \cdots + a_0}}
}
\begin{itemize}
  \item \textbf{P\^oles} $p_i$ : racines du d\'enom. $\Rightarrow$ modes dynamiques
  \item \textbf{Z\'eros} $z_i$ : racines du num\'erateur
  \item \textbf{Ordre} $n$ = degr\'e d\'enom. \quad \textbf{Degr\'e rel.} $d = n - m$
\end{itemize}
\end{tcolorbox}

\begin{tcolorbox}[colback=ch2orange!2,colframe=ch2orange!80!black,title={Trois formes de $G(s)$}]
\textbf{1. Forme de Bode} (normalis\'ee) :
\tightmath{
G(s) = \frac{K}{s^\alpha}\cdot\frac{(1{+}\tau_{z_1}s)(1{+}2\zeta_z\frac{s}{\omega_z}{+}\frac{s^2}{\omega_z^2})\cdots}{(1{+}\tau_{p_1}s)(1{+}2\zeta_p\frac{s}{\omega_p}{+}\frac{s^2}{\omega_p^2})\cdots}
}
\astuce{C'est la forme \`a utiliser pour tracer le Bode !}

\textbf{2. Forme factoris\'ee} (p\^oles-z\'eros) :
\tightmath{
G(s) = \frac{b_m}{a_n}\cdot\frac{(s{-}z_1)(s{-}z_2)\cdots}{(s{-}p_1)(s{-}p_2)\cdots}
}
\textbf{3. Forme polynomiale} : num/d\'enom d\'evelopp\'es.

\vars{\textbf{Passage} : $K = b_0/a_0$ (forme Bode, $\alpha{=}0$). Gain permanent : $K = \lim_{s\to 0} s^\alpha G(s)$. Constantes : $\tau_i = -1/p_i$.}
\end{tcolorbox}

\begin{tcolorbox}[colback=ch2orange!4,colframe=ch2orange!80!black,title={Changement de forme}]
\begin{minipage}[t]{0.48\linewidth}
\textbf{Laplace $\rightarrow$ Bode}

\tightmath{
\frac{K}{s^2 + 2\sigma s + \omega_n^2}
= \frac{K/\omega_n^2}{1 + \frac{2\zeta}{\omega_n}s + \frac{s^2}{\omega_n^2}}
}
\tightmath{
\omega_n = \sqrt{\sigma^2 + \omega_0^2}
\qquad
\zeta = \frac{\sigma}{\omega_n}
}
\vars{Aussi : $(s+a)=a\!\left(1+\frac{s}{a}\right)$, donc $\tau=\frac{1}{a}$ et $\omega_c=a$.}
\end{minipage}\hfill
\begin{minipage}[t]{0.48\linewidth}
\textbf{Bode $\rightarrow$ Laplace}

\begin{center}
\resizebox{0.98\linewidth}{!}{$\displaystyle
\frac{K}{1 + \frac{2\zeta s}{\omega_n} + \frac{s^2}{\omega_n^2}}
= \frac{K\omega_n^2}{s^2 + 2\zeta\omega_n s + \omega_n^2}
$}
\end{center}
\tightmath{
\sigma = \zeta\omega_n
\qquad
\omega_0 = \omega_n\sqrt{1-\zeta^2}
}
\tightmath{
K_{\text{Laplace}} = K_{\text{Bode}}\,\omega_n^2
}
\end{minipage}
\vars{Id\'ee : en forme de Bode, la constante dans les parenth\`eses vaut $1$ ; en forme de Laplace, on red\'eveloppe et on r\'einjecte les facteurs sortis.}
\end{tcolorbox}

\begin{tcolorbox}[colback=ch2orange!4,colframe=ch2orange!80!black,title={Type $\alpha$, Retard pur, Gain permanent}]
\textbf{Type $\alpha$} = nb de p\^oles en $s{=}0$ (int\'egrateurs en BO) :
\begin{itemize}
  \item $\alpha{=}0$ : erreur statique $\neq 0$ pour \'echelon
  \item $\alpha{=}1$ : $1/s$, \'elimine erreur pour \'echelon
  \item $\alpha{=}2$ : $1/s^2$, \'elimine erreur pour rampe
\end{itemize}
\textbf{Gain permanent} : $\boxed{K = \lim_{s\to 0} s^\alpha \cdot G(s)}$

\textbf{Retard pur} : $G_{ret}(s) = e^{-sT_r}$
\begin{itemize}
  \item Module : $|e^{-\jw T_r}| = 1$ (inchang\'e !)
  \item Phase : $\angle e^{-\jw T_r} = -\omega T_r$ rad $= -\omega T_r \cdot \frac{180°}{\pi}$
\end{itemize}
\piege{Le retard fait chuter la phase ind\'efiniment $\Rightarrow$ d\'estabilisant !}
\end{tcolorbox}

\begin{tcolorbox}[colback=ch2orange!2,colframe=ch2orange!80!black,title={Alg\`ebre des sch\'emas blocs}]
\tightmath{
\textbf{S\'erie} :\ G = G_1 \cdot G_2 \cdot \ldots \cdot G_n
}
\tightmath{
\textbf{Parall\`ele} :\ G = G_1 \pm G_2
}
\tightmath{
\textbf{BF (r\'etro. $-$)} :\ \boxed{G_{BF} = \frac{G_1}{1 + G_1 G_2}} \quad \textbf{(r\'etro. $+$)} :\ \frac{G_1}{1 - G_1 G_2}
}
\piege{Signe $+$ au d\'enom. pour r\'etroaction \textbf{n\'egative}. Signe $-$ si positive.}

\begin{center}
\renewcommand{\arraystretch}{1}
\scriptsize
\begin{tabular}{@{}ll@{}}
\toprule
Op\'eration & R\'esultat \\
\midrule
Sommateur avant bloc $G$ & Ajouter $1/G$ dans branche \\
Sommateur apr\`es bloc $G$ & Ajouter $G$ dans branche \\
Pt pr\'el\`evement avant $G$ & Ajouter $G$ dans branche \\
Pt pr\'el\`evement apr\`es $G$ & Ajouter $1/G$ dans branche \\
\bottomrule
\end{tabular}
\end{center}
\end{tcolorbox}

\begin{tcolorbox}[colback=ch2orange!4,colframe=ch2orange!80!black,title={Cha\^ine d'action \& Mason}]
\textbf{Cha\^ine d'action} : $G = G_1G_2\cdots G_n$. \quad
\textbf{Boucle ouverte} : $L(s)=G_rG_aG_pG_m$.

\tightmath{
\boxed{G_{BF}(s)=\frac{\text{cha\^ine directe}}{1 \pm \text{gain de boucle}}}
}
\tightmath{
\text{Retour }(-):\ \boxed{G_{BF}=\frac{G}{1+GH}}
\text{  Retour }(+):\ \boxed{G_{BF}=\frac{G}{1-GH}}
}
\begin{center}
\scriptsize
\resizebox{0.98\linewidth}{!}{%
\begin{tikzpicture}[
  >=Latex,
  node distance=4mm and 5mm,
  blk/.style={draw,minimum width=9mm,minimum height=4.5mm,inner sep=1pt},
  sum/.style={draw,circle,inner sep=0pt,minimum size=4.2mm}
]
  \node (in1) {$x$};
  \node[blk,right=of in1] (g1) {$G_1$};
  \node[blk,right=of g1] (g2) {$G_2$};
  \node[right=of g2] (out1) {$y$};
  \draw[->] (in1) -- (g1);
  \draw[->] (g1) -- (g2);
  \draw[->] (g2) -- (out1);
  \node[below=1mm of g1.south east] {$P_1 = G_1G_2$};

  \begin{scope}[xshift=40mm]
    \node (w) {$w$};
    \node[sum,right=of w] (s1) {$\scriptstyle\sum$};
    \node[blk,right=of s1] (g) {$G$};
    \node[right=of g] (y1) {$y$};
    \node[blk,below=7mm of g] (h1) {$H$};
    \draw[->] (w) -- (s1);
    \draw[->] (s1) -- (g);
    \draw[->] (g) -- (y1);
    \draw[->] (y1) |- (h1);
    \draw[->] (h1) -| (s1);
    \node[above=0.4mm of s1] {\tiny $+$};
    \node[below=0.4mm of s1, xshift=2mm] {\tiny $-$};
    \node[below=1mm of h1] {$\text{boucle}=-GH \Rightarrow \Delta=1+GH$};
  \end{scope}

  \begin{scope}[xshift=82mm]
    \node (w2) {$w$};
    \node[sum,right=of w2] (s2) {$\scriptstyle\sum$};
    \node[blk,right=of s2] (gp) {$G$};
    \node[right=of gp] (y2) {$y$};
    \node[blk,below=7mm of gp] (h2) {$H$};
    \draw[->] (w2) -- (s2);
    \draw[->] (s2) -- (gp);
    \draw[->] (gp) -- (y2);
    \draw[->] (y2) |- (h2);
    \draw[->] (h2) -| (s2);
    \node[above=0.4mm of s2] {\tiny $+$};
    \node[below=0.4mm of s2, xshift=2mm] {\tiny $+$};
    \node[below=1mm of h2] {$\text{boucle}=+GH \Rightarrow \Delta=1-GH$};
  \end{scope}
\end{tikzpicture}
}
\end{center}
\vars{En BF unitaire : $L(s)=G_{BO}(s)$. En cas de doute, repartir de $y=G(u\mp Hy)$.}
\vars{Mason : $T=\dfrac{P_1}{\Delta}$ avec $\Delta = 1 - \sum(\text{gains de boucle}) + \cdots$ ; donc une boucle n\'egative de gain $-GH$ donne bien $1+GH$.}
\end{tcolorbox}

\begin{tcolorbox}[colback=ch2orange!4,colframe=ch2orange!80!black,title={D\'ecomposition en fractions partielles}]
\textbf{P\^oles simples} $p_i$ :
\tightmath{
\frac{N(s)}{(s{-}p_1)(s{-}p_2)\cdots} = \sum_i \frac{A_i}{s - p_i} \quad \boxed{A_i = \lim_{s\to p_i}(s{-}p_i) G(s)}
}
\textbf{P\^ole double} en $p$ :
\tightmath{
\frac{A}{(s{-}p)^2} + \frac{B}{s{-}p} + \cdots
}
$A = \lim_{s\to p}(s{-}p)^2 G(s)$, \; $B = \lim_{s\to p}\frac{d}{ds}[(s{-}p)^2 G(s)]$

\textbf{P\^oles complexes} $-\sigma \pm j\omega$ : garder $\frac{As+B}{(s{+}\sigma)^2+\omega^2}$ et utiliser table.

\vars{Conversions : Hz$\to$rad/s : $\omega{=}2\pi f$. dB$\to$lin : $|G|{=}10^{G_{\text{dB}}/20}$. lin$\to$dB : $20\log_{10}|G|$. $6$~dB $\approx \times 2$. $20$~dB $= \times 10$.}
\end{tcolorbox}

\begin{tcolorbox}[colback=ch2orange!2,colframe=ch2orange!80!black,title={Inversion de Laplace : recette d'examen}]
\begin{enumerate}
  \item V\'erifier si la fraction est propre (sinon division polynomiale)
  \item Factoriser le d\'enominateur en p\^oles simples/r\'ep\'et\'es/complexes
  \item Choisir la bonne forme : $\frac{A}{s-p}$, $\frac{A}{(s-p)^2}$, $\frac{As+B}{(s+\sigma)^2+\omega^2}$
  \item Identifier les coefficients (m\'ethode des limites ou identification)
  \item Appliquer la table de Laplace, puis simplifier l'expression temporelle
\end{enumerate}
\tightmath{
\frac{5}{s(s+2)} = \frac{2.5}{s} - \frac{2.5}{s+2} \Rightarrow y(t)=2.5\big(1-e^{-2t}\big)
}
\astuce{Pour un \'echelon : $Y(s)=G(s)\,\frac{A}{s}$, le p\^ole suppl\'ementaire en $s=0$ est normal.}
\end{tcolorbox}

\begin{tcolorbox}[colback=ch2orange!4,colframe=ch2orange!80!black,title={M\'ethodologie : esquisse de la r\'ep. temporelle},top=0.3mm,bottom=0.3mm,before skip=0.8pt,after skip=0.8pt]
\textbf{\'Etapes pour esquisser $y(t)$ \`a un signal d'entr\'ee quelconque :}
\begin{itemize}
  \item Trouver $K$, $y_\infty$ et le retard $T_r$
  \item P\^oles r\'eels $\Rightarrow$ ap\'eriodique ; complexes $\Rightarrow$ oscillatoire
  \item Calculer $T_{\text{r\'eg}} = 3/|\sigma|$ et $T_0 = 2\pi/\omega_0$
  \item D\'ecaler de $T_r$ et sommer les r\'eponses aux \'echelons successifs
\end{itemize}
\vars{Exemple : $G(s)=\dfrac{10}{(s+2)^2+16}e^{-s}$, p\^oles $-2\pm j4$, $y_\infty=0.4$, $T_{\text{r\'eg}}=1.5$ s, $T_0\approx1.57$ s, $T_r=1$ s.}
\begin{center}
\begin{tikzpicture}[scale=0.47]
 
% ---- Paramètres de la courbe ----
% G(s) = 10/((s+2)^2+16) * e^{-s}
% Réponse simulée manuellement avec retard Tr=1
% pôles: -2 ± j4, omega_d=4, sigma=2, y_inf = 0.4
 
\def\Tr{1.0}     % retard
\def\Treg{2.4}   % temps de réglage mesuré depuis Tr
\def\Tzero{1.6}  % pseudo-période T0
 
% ---- Couleurs ----
\definecolor{courbe}{RGB}{160, 82, 10}
\definecolor{orangetr}{RGB}{220,140,40}
\definecolor{bleutreg}{RGB}{50,100,200}
\definecolor{rougeto}{RGB}{200,40,40}
\definecolor{grisaxe}{RGB}{80,80,80}
\definecolor{bandefill}{RGB}{200,230,210}

% ---- Point lu sur la courbe pour Treg ----
\def\yTreg{2.05}

% ============================================================
% Axes (dessinés d'abord pour rester en arrière-plan)
% ============================================================
\draw[-{Stealth[length=7pt]}, grisaxe, thick]
  (-0.2, 0) -- (9.8, 0) node[right, font=\small] {$t$};
\draw[-{Stealth[length=7pt]}, grisaxe, thick]
  (0, -0.8) -- (0, 4.4) node[above, font=\small] {$y$};
 
% ---- Bande ±5% ----
\fill[bandefill, opacity=0.45]
  (0, 2.05) rectangle (9.5, 2.55);
\draw[gray!50, dashed, thin] (0, 2.05) -- (9.5, 2.05);
\draw[gray!50, dashed, thin] (0, 2.30) -- (9.5, 2.30);  % y_inf
\draw[gray!50, dashed, thin] (0, 2.55) -- (9.5, 2.55);
 
% ---- Consigne w(t) ----
\draw[grisaxe, thick, magenta] (0,0) -- (0, 1.5) -- (9.5, 1.5)
  node[right, font=\small, magenta] {$w(t)$};
 
% ---- Courbe de réponse y(t) ----
% Retard: plat à 0 jusqu'à t=Tr
% Puis réponse oscillante amortie
\draw[courbe, very thick, line join=round]
  (0, 0) -- (\Tr, 0)
  .. controls (\Tr+0.15, 0) and (\Tr+0.4, 0.5) ..
  (\Tr+0.7, 2.0)
  .. controls (\Tr+0.85, 2.8) and (\Tr+1.0, 3.5) ..
  (\Tr+1.2, 3.6)
  .. controls (\Tr+1.5, 3.65) and (\Tr+1.7, 3.5) ..
  (\Tr+2.0, 2.7)
  .. controls (\Tr+2.3, 1.9) and (\Tr+2.5, 1.6) ..
  (\Tr+2.8, 1.85)
  .. controls (\Tr+3.1, 2.1) and (\Tr+3.3, 2.35) ..
  (\Tr+3.5, 2.42)
  .. controls (\Tr+3.7, 2.48) and (\Tr+3.9, 2.38) ..
  (\Tr+4.2, 2.28)
  .. controls (\Tr+4.5, 2.20) and (\Tr+4.8, 2.22) ..
  (\Tr+5.0, 2.26)
  .. controls (\Tr+5.2, 2.30) and (\Tr+5.5, 2.32) ..
  (\Tr+6.0, 2.30)
  -- (9.5, 2.30);
 
% ---- y_inf label ----
\node[right, font=\small\itshape, text=grisaxe] at (9.5, 2.30) {$y_\infty$};
 
% ============================================================
% ANNOTATION 1 : Tr (retard) — flèche orange en bas
% ============================================================
\draw[orangetr, dashed, thin] (\Tr, 0) -- (\Tr, -0.55);
\draw[orangetr,
  {Stealth[length=5pt]}-{Stealth[length=5pt]}]
  (0, -0.45) -- (\Tr, -0.45);
\fill[courbe] (\Tr, 0) circle (1.8pt);
\node[below, font=\small\bfseries, text=orangetr] at (\Tr/2, -0.55)
  {$T_r$};
 
% ============================================================
% ANNOTATION 2 : Treg — double flèche bleue depuis Tr
% ============================================================
\def\Tregend{4.025}  % Tr + Treg sur l'axe x
\draw[bleutreg, dashed, thin] (\Tregend, -0.55) -- (\Tregend, \yTreg);
\draw[bleutreg,
  {Stealth[length=5pt]}-{Stealth[length=5pt]}]
  (\Tr, -0.45) -- (\Tregend, -0.45);
\fill[bleutreg] (\Tregend, \yTreg) circle (3pt);
\node[below, font=\small\bfseries, text=bleutreg] at ({(\Tr+\Tregend)/2}, -0.45)
  {$T_{\mathrm{reg}}$};
 
% ============================================================
% ANNOTATION 3 : T0 (pseudo-période) — flèche rouge en haut
% ============================================================
% T0 va du premier pic au second pic
% Premier pic approx à x = Tr+1.2, second pic à Tr+1.2+T0
\def\xpic{2.3}       % premier pic
\def\xpicsec{4.58}   % second "pic" (correspond à T0=1.57)
\def\ypic{3.75}
 
\draw[rougeto, dashed, thin] (\xpic, 3.65) -- (\xpic, \ypic+0.1);
\draw[rougeto, dashed, thin] (\xpicsec, 2.47) -- (\xpicsec, \ypic+0.1);
\draw[rougeto,
  {Stealth[length=5pt]}-{Stealth[length=5pt]}]
  (\xpic, \ypic) -- (\xpicsec, \ypic);
\node[above, font=\small\bfseries, text=rougeto] at ({(\xpic+\xpicsec)/2}, \ypic)
  {$T_0$};
 
\end{tikzpicture}
\end{center}
\end{tcolorbox}

\begin{tcolorbox}[colback=ch2orange!2,colframe=ch2orange!80!black,title={Identification d'un syst\`eme oscillant}]
\tightmath{
G(s)=\frac{K_0}{1+\frac{2\zeta}{\omega_n}s+\frac{s^2}{\omega_n^2}}
}
\tightmath{
K_0=\frac{y(\infty)}{\Delta U}
\qquad
\omega_0=\omega_n\sqrt{1-\zeta^2}
}
\tightmath{
\boxed{\zeta=\frac{|\ln D_1|}{\sqrt{\pi^2+\ln^2 D_1}}}
\qquad
D_1=\frac{y_{\max}-y(\infty)}{y(\infty)}
}
\tightmath{
\boxed{\omega_n=\frac{2\pi}{T_p\sqrt{1-\zeta^2}}}
\qquad
T_p=\text{p\'eriode des oscillations}
}
\vars{$D_1$ = 1\textsuperscript{er} d\'epassement relatif. Pour un \'echelon de hauteur $\Delta U$, on a $y(\infty)=K_0\Delta U$.}
\end{tcolorbox}

\begin{tcolorbox}[colback=ch2orange!4,colframe=ch2orange!80!black,title={Lecture d'un diagramme de Bode}]
\textbf{Type $\alpha$} $\leftarrow$ pente \`a basse fr\'equence :
\tightmath{
  \alpha = -\frac{\text{pente BF (dB/d\'ec)}}{20}
}
\begin{center}
\renewcommand{\arraystretch}{1.0}
\scriptsize
\begin{tabular}{@{}ccc@{}}
\toprule
Pente BF & Type $\alpha$ & Exemple \\
\midrule
$0$ dB/d\'ec & $0$ & $G(s)=\frac{K}{1+\tau s}$ \\
$-20$ dB/d\'ec & $1$ & $G(s)=\frac{K}{s(1+\tau s)}$ \\
$-40$ dB/d\'ec & $2$ & $G(s)=\frac{K}{s^2(\cdots)}$ \\
\bottomrule
\end{tabular}
\end{center}

\textbf{Degr\'e relatif $d$} $\leftarrow$ pente \`a haute fr\'equence :
\tightmath{
  d = \frac{-\text{pente HF (dB/d\'ec)}}{20}
}
\begin{center}
\begin{tikzpicture}[>=Stealth, scale=0.60, transform shape,
  lbl/.style={font=\fontsize{11}{12.5}\selectfont}]
\colorlet{asym}{red!75!black}
\colorlet{bode}{ch2orange!80!black}
\colorlet{zcol}{green!40!black}
\colorlet{pcol}{red!60!black}

% Grille legere
\foreach \x in {2.5,5.0} \draw[gray!22, thin] (\x,-0.3) -- (\x,5.7);

% Axes
\draw[->, thick] (-0.3,0) -- (7.8,0)
  node[right, lbl]{$\omega$ (log)};
\draw[->, thick] (0,-0.4) -- (0,5.9)
  node[above, lbl]{$|G|_{\mathrm{dB}}$};

% Asymptote pointillee rouge
% alpha=1 : pente -20 dB/dec jusqu'a wz
% zero a wz : pente 0 jusqu'a wp
% double pole a wp : pente -40 dB/dec -> d=2
\draw[asym, dashed, line width=0.9pt]
  (0,5.0) -- (2.5,2.5) -- (5.0,2.5) -- (6.25,0);

% Courbe Bode lisse (approximation avec pic de resonance a wp)
\draw[bode, line width=1.2pt]
  (0,4.85)
  .. controls (0.8,4.10) and (1.7,3.20) ..
  (2.4,2.68)
  .. controls (2.90,2.48) and (3.80,2.50) ..
  (4.55,2.58)
  .. controls (4.80,2.72) and (4.93,3.20) ..
  (5.05,3.55)
  .. controls (5.18,3.25) and (5.38,2.45) ..
  (5.55,1.62)
  .. controls (5.78,0.92) and (6.02,0.38) ..
  (6.28,0.06);

% Pointilles verticaux aux frequences cles
\draw[gray!45, dashed, thin] (2.5,0) -- (2.5,2.5);
\draw[gray!45, dashed, thin] (5.05,0) -- (5.05,3.55);

% Marqueurs sur la courbe
\fill[zcol] (2.5,2.5)  circle (2.5pt);
\fill[pcol] (5.05,3.55) circle (2.5pt);

% Labels frequences
\node[below, lbl] at (2.5,-0.05)  {$\omega_z$};
\node[below, lbl] at (5.05,-0.05) {$\omega_p$};

% Annotation pente BF -> alpha
\node[asym, lbl, above right] at (0.1,4.55) {$-20$\,dB/d\'ec};
\node[lbl, above right]       at (0.8,4.2) {$\Rightarrow\boldsymbol{\alpha=1}$};

% Annotation zero
\draw[zcol, ->, thin] (2.5,2) -- (2.5,2.4);
\node[zcol, lbl, below] at (2.5,2) {z\'ero\;($+20$\,dB/d\'ec)};

% Annotation pole(s)
\draw[pcol, ->, thin] (5.62,4.15) -- (5.23,3.70);
\node[pcol, lbl, right] at (3.85,4.38) {p\^ole(s) ($-20$\,dB/d\'ec chacun)};
\node[pcol, lbl, right] at (8.75,4.38) {+ r\'esonance si $\zeta{<}0.7$};

% Annotation pente HF -> d
\node[asym, lbl, right] at (5.80,1.95) {$-40$\,dB/d\'ec};
\node[lbl, right]        at (5.80,1.32) {$\Rightarrow\boldsymbol{d=2}$};

\end{tikzpicture}
\end{center}
\piege{\textbf{Sur graph} : lire la pente avant tout p\^ole/z\'ero (zone plate ou rampe), pas autour d'un pic de r\'esonance.}
\end{tcolorbox}

\begin{tcolorbox}[colback=ch2orange!2,colframe=ch2orange!80!black,title={Astuces \& Pi\`eges Ch.2},top=0.4mm,bottom=0.4mm,before skip=1pt,after skip=1pt]
\begin{itemize}
  \item VF : $\lim_{t\to\infty}f(t)=\lim_{s\to0}sF(s)$ \textbf{ssi stable}
  \item Bode : normaliser, ex. $\frac{2}{s+4}=\frac{0.5}{1+0.25s}$
  \item BF : r\'etroaction $(-)\Rightarrow 1+G_1G_2$, r\'etroaction $(+)\Rightarrow 1-G_1G_2$
  \item Retard pur : $|e^{-sT_r}|=1$, p\^oles inchang\'es
  \item \'Eq. diff $\to G(s)$ : Laplace CI nulles puis isoler $Y/U$
\end{itemize}
\end{tcolorbox}
